% !TEX root = ../main.tex

\chapter{بحث و تحلیل نتایج}

% =========================================================
% 7-1
% =========================================================
\section{مقدمه}

در فصل قبل، نتایج کمی مربوط به ارزیابی مدل‌های \lr{ASR}، \lr{LLM} و
\lr{TTS} ارائه شدند.
در این فصل، این نتایج از دیدگاه طراحی سامانه تحلیل می‌شوند.

هدف اصلی، بررسی این موضوع است که چرا برخی مدل‌ها عملکرد بهتری
داشته‌اند و چگونه باید میان دقت، سرعت و مصرف منابع توازن
برقرار شود.

همچنین محدودیت‌های آزمایش‌ها و عواملی که می‌توانند بر
تفسیر نتایج اثر بگذارند بررسی خواهند شد.


% =========================================================
% 7-2
% =========================================================
\section{تحلیل نتایج تشخیص گفتار}

در آزمایش \lr{ASR}، مدل \lr{FastConformer Hybrid} کمترین نرخ خطای واژه
را در میان مدل‌های بررسی‌شده ارائه کرد.

مقدار \lr{WER} این مدل حدود:

\begin{equation}
WER = 0.1455
\end{equation}

بود، در حالی که برای \lr{wav2vec2}، \lr{Vosk} و \lr{Whisper} به‌ترتیب
مقادیر بالاتری مشاهده شد.

این نتیجه نشان می‌دهد که \lr{FastConformer} علاوه بر سرعت مناسب،
سازگاری خوبی با داده فارسی مورد استفاده در این آزمایش داشته است.

یکی از دلایل احتمالی این عملکرد، آموزش و تنظیم اختصاصی‌تر
مدل برای زبان فارسی و استفاده از معماری \lr{FastConformer} همراه
با ساختار Hybrid CTC/RNN-T است.

با این حال، نمی‌توان نتیجه گرفت که این مدل در تمام مجموعه‌های
فارسی یا تمام شرایط محیطی الزاماً بهترین خواهد بود.
نوع داده، نحوه نرمال‌سازی متن و شرایط صوتی می‌توانند بر \lr{WER}
اثر قابل توجهی داشته باشند.


% ---------------------------------------------------------
\subsection{تحلیل عملکرد \lr{wav2vec2}}

مدل \lr{wav2vec2} از نظر سرعت بهترین نتیجه را ارائه کرد.
\lr{RTF} بسیار پایین آن نشان می‌دهد که مدل می‌تواند صوت را با سرعت
بسیار بیشتر از زمان واقعی پردازش کند.

با وجود این، \lr{WER} آن به‌طور قابل توجهی بیشتر از \lr{FastConformer}
بود.

بنابراین \lr{wav2vec2} نمونه‌ای از یک توازن میان سرعت و دقت است.
اگر سرعت مهم‌ترین عامل باشد، این مدل گزینه مناسبی است، اما
برای پروژه حاضر دقت تشخیص نیز اهمیت زیادی دارد؛ زیرا خطای \lr{ASR}
مستقیماً به مدل زبانی منتقل می‌شود.


% ---------------------------------------------------------
\subsection{تحلیل عملکرد \lr{Whisper}}

\lr{Whisper Large-v3} در این بنچمارک بالاترین مصرف حافظه و
همچنین \lr{WER} بالاتری نسبت به سایر مدل‌های بررسی‌شده داشت.

این نتیجه در نگاه اول ممکن است غیرمنتظره باشد، زیرا \lr{Whisper}
یک مدل بزرگ و چندزبانه قدرتمند است.

با این حال، اندازه مدل به‌تنهایی تضمین‌کننده عملکرد بهتر در
یک زبان و مجموعه‌داده خاص نیست.

عواملی مانند:

\begin{itemize}
    \item نوع داده آموزشی؛
    \item تفاوت میان گفتار فارسی آزمون و داده‌های آموزش؛
    \item روش قطعه‌بندی صوت؛
    \item تنظیمات رمزگشایی؛
    \item و نرمال‌سازی متن
\end{itemize}

می‌توانند در نتیجه تأثیر داشته باشند.

همچنین معماری خودبازگشتی \lr{Whisper} هزینه زمانی بیشتری نسبت به
مدل‌هایی مانند \lr{wav2vec2} ایجاد می‌کند.

در نتیجه، برای هدف اجرای محلی، هزینه محاسباتی بالاتر \lr{Whisper}
در این آزمایش با بهبود دقت جبران نشده است.


% ---------------------------------------------------------
\subsection{تحلیل عملکرد \lr{Vosk}}

مهم‌ترین مزیت \lr{Vosk} در این آزمایش عدم نیاز به \lr{GPU} بود.

این ویژگی می‌تواند برای سیستم‌هایی که پردازنده گرافیکی ندارند
مهم باشد.

با این حال، سرعت و دقت آن نسبت به \lr{FastConformer} در این
آزمایش ضعیف‌تر بود.

بنابراین \lr{Vosk} بیشتر به‌عنوان یک گزینه کم‌منبع قابل توجه است تا
مدلی که بهترین کیفیت کلی را ارائه کند.


% =========================================================
% 7-3
% =========================================================
\section{تحلیل نتایج مدل‌های زبانی}

نتایج کنترل‌شده وظیفه چندگزینه‌ای نشان داد که \lr{Gemma 2 9B}
با دقت کل \lr{46.86\%} بهترین عملکرد را دارد. برتری این مدل در
دانش عمومی و ادبیات نشان می‌دهد که ظرفیت بیشتر آن در بازیابی و
به‌کارگیری دانش فارسی سودمند بوده است. با این حال، برتری
\lr{Qwen2.5 7B} در ریاضی و منطق با دقت \lr{42.57\%} نشان می‌دهد
که رتبه مدل‌ها به نوع وظیفه وابسته است و یک مدل در همه حوزه‌ها
لزوماً بهترین نیست.

فاصله قابل توجه \lr{Gemma 2 9B} و \lr{Gemma 2 2B} در دقت کل، یعنی
\lr{46.86\%} در برابر \lr{25.71\%}، نشان می‌دهد که نتیجه آزمایش
مقدماتی \lr{FarsInstruct} را نباید به‌عنوان برتری عمومی مدل کوچک‌تر
تفسیر کرد. ترکیب تصادفی وظایف در آن آزمایش، مقایسه حوزه‌به‌حوزه را
دشوار می‌کرد؛ ازاین‌رو رتبه‌بندی اصلی بر نتایج \lr{ParsiNLU} استوار است.


% ---------------------------------------------------------
\subsection{تفسیر معیارهای درک مطلب}

در ارزیابی \lr{ParsBench RC}، تطابق دقیق معیار سخت‌گیرانه‌ای است و
تنها در صورت یکسان بودن پاسخ پیش‌بینی‌شده و پاسخ مرجع امتیاز کامل
می‌دهد. \lr{F1} توکنی به همپوشانی اجزای پاسخ حساس است، در حالی که
شباهت معنایی پاسخ‌های هم‌معنا با صورت‌بندی متفاوت را بهتر پوشش
می‌دهد. بنابراین تفسیر نتایج درک مطلب باید بر ترکیب این سه معیار
استوار باشد، نه یک معیار منفرد.

برای مثال:

\begin{center}
«تهران پایتخت ایران است.»
\end{center}

و:

\begin{center}
«پایتخت کشور ایران، تهران است.»
\end{center}

معنای بسیار مشابهی دارند، اما تطابق واژگانی آن‌ها کامل نیست.

به همین دلیل، گزارش هم‌زمان تطابق دقیق، \lr{F1} توکنی و شباهت
معنایی تصویر کامل‌تری از کیفیت پاسخ ارائه می‌کند. با این حال،
ارزیابی انسانی در آینده می‌تواند اعتبار نتایج را بیشتر کند.


% ---------------------------------------------------------
\subsection{توازن اندازه مدل و کیفیت}

اجرای همه مدل‌ها با کوانتیزه‌سازی چهاربیتی \lr{NF4} هزینه حافظه را
کاهش داد و امکان مقایسه مدل‌های \lr{2B} تا \lr{9B} را در یک چارچوب
مشترک فراهم کرد. با وجود این، مدل‌های بزرگ‌تر همچنان به منابع
بیشتری نیاز دارند. در نتیجه، \lr{Gemma 2 9B} انتخاب مبتنی بر بیشترین
دقت کلی است؛ اما در سخت‌افزار محدودتر، انتخاب یک مدل کوچک‌تر می‌تواند
به بهای کاهش دقت، تأخیر و مصرف حافظه را کاهش دهد. همچنین اگر وظیفه
اصلی سامانه بر استدلال ریاضی و منطقی متمرکز باشد، \lr{Qwen2.5 7B}
گزینه شایان توجهی است.


% =========================================================
% 7-4
% =========================================================
\section{تحلیل نتایج تبدیل متن به گفتار}

نتایج \lr{TTS} یک توازن واضح میان کیفیت و سرعت را نشان دادند.

مدل \lr{Piper} بالاترین امتیاز \lr{NISQA} و کمترین \lr{WER} و \lr{CER} را
ارائه کرد.
بنابراین خروجی آن از نظر کیفیت ادراکی و قابل‌فهم‌بودن در
آزمایش انجام‌شده بهترین عملکرد را داشت.

با این حال، \lr{RTF} آن نسبت به دو مدل دیگر بیشتر بود.

در مقابل، \lr{MMS VITS} سریع‌ترین مدل بود، اما \lr{WER} و \lr{CER} بالاتر
و کیفیت ادراکی پایین‌تری نسبت به \lr{Piper} داشت.

این نتیجه نشان می‌دهد که انتخاب مدل \lr{TTS} برای یک چت‌بات
تنها بر اساس سرعت کافی نیست.


% ---------------------------------------------------------
\subsection{انتخاب \lr{Piper}}

با وجود سرعت کمتر \lr{Piper}، مقدار \lr{RTF} آن همچنان کمتر از یک بود:

\begin{equation}
RTF \approx 0.192
\end{equation}

بنابراین مدل همچنان سریع‌تر از زمان واقعی صوت تولید می‌کند.

از آنجا که کیفیت صوت مستقیماً برای کاربر قابل درک است،
بهبود قابل توجه کیفیت \lr{Piper} می‌تواند افزایش محدود زمان تولید
را توجیه کند.

به همین دلیل، \lr{Piper} به‌عنوان کاندید اصلی \lr{TTS} انتخاب شده است.


% =========================================================
% 7-5
% =========================================================
\section{توازن میان کیفیت و سرعت در سامانه نهایی}

نتایج سه بخش نشان می‌دهند که انتخاب مدل تنها بر اساس یک معیار
منفرد مناسب نیست.

برای مثال:

\begin{itemize}

    \item \lr{wav2vec2} سریع‌ترین \lr{ASR} بود، اما دقیق‌ترین نبود؛

    \item \lr{FastConformer} دقیق‌ترین \lr{ASR} بود و همچنان سرعت
    مناسبی داشت؛

    \item \lr{Gemma 2 9B} بیشترین دقت کل را در پرسش‌های چندگزینه‌ای
    به دست آورد و \lr{Qwen2.5 7B} در ریاضی و منطق بهترین بود؛

    \item و \lr{Piper} بهترین کیفیت \lr{TTS} را با هزینه سرعت بیشتر
    ارائه کرد.

\end{itemize}

بنابراین ترکیب پیشنهادی:

\begin{center}
\lr{FastConformer}
$\longrightarrow$
\lr{Gemma 2 2B}
$\longrightarrow$
\lr{Piper}
\end{center}

به‌عنوان گزینه کم‌منبع سامانه پیاده‌سازی‌شده در نظر گرفته شده است.
در کاربردی که دقت وظایف فارسی اولویت بیشتری از محدودیت منابع دارد،
\lr{Gemma 2 9B} جایگزین مناسب‌تری است. تصمیم نهایی میان این دو باید
با اندازه‌گیری هم‌زمان دقت، تأخیر و حافظه در آزمون انتهابه‌انتها
انجام شود.


% =========================================================
% 7-6
% =========================================================
\section{انتشار خطا میان مؤلفه‌ها}

یکی از محدودیت‌های معماری زنجیره‌ای، انتقال خطا از یک مرحله
به مرحله بعد است.

اگر \lr{ASR} بخشی از درخواست را اشتباه تشخیص دهد، مدل زبانی
پاسخ خود را بر اساس همان متن اشتباه تولید خواهد کرد.

در نتیجه، کیفیت سامانه نهایی را نمی‌توان تنها با جمع‌کردن
عملکرد مستقل مؤلفه‌ها پیش‌بینی کرد.

این موضوع ضرورت ارزیابی انتهابه‌انتهای سامانه را نشان می‌دهد.

به‌طور مشابه، ممکن است مدل زبانی پاسخ صحیحی تولید کند اما
\lr{TTS} بخشی از متن را به‌صورت نامفهوم تلفظ کند.

بنابراین کیفیت نهایی تابع تعامل تمام مراحل است.


% =========================================================
% 7-7
% =========================================================
\section{محدودیت‌های ارزیابی}

نتایج ارائه‌شده باید با توجه به چند محدودیت تفسیر شوند.

نخست، عملکرد مدل‌ها به مجموعه‌داده استفاده‌شده وابسته است و
ممکن است روی داده‌های دیگر نتایج متفاوتی حاصل شود.

دوم، معیارهای خودکار کیفیت مدل زبانی و \lr{TTS} جایگزین کامل
ارزیابی انسانی نیستند.

سوم، شرایط نویزی و گفتار کاملاً طبیعی هنوز باید به‌صورت
گسترده‌تر آزمایش شوند.

چهارم، ارزیابی کامل VAD، مدیریت بافت مکالمه و تأخیر
انتهابه‌انتها هنوز باید تکمیل شود.

در نتیجه، نتایج فعلی بیشتر برای انتخاب کاندیدهای مناسب و
طراحی سامانه استفاده می‌شوند و نباید به‌عنوان یک رتبه‌بندی
عمومی برای تمام کاربردهای فارسی تفسیر شوند.


% =========================================================
% 7-8
% =========================================================
\section{جمع‌بندی}

در این فصل نتایج آزمایش‌ها از دیدگاه طراحی سامانه تحلیل شدند.

\lr{FastConformer} بهترین توازن دقت و سرعت را در بخش \lr{ASR} ارائه کرد.
\lr{Gemma 2 9B} در ارزیابی چندگزینه‌ای \lr{ParsiNLU} بیشترین دقت کل
را نشان داد و اختلاف عملکرد میان دسته‌ها، اهمیت انتخاب مدل بر اساس
نوع وظیفه را آشکار کرد.
در بخش \lr{TTS} نیز \lr{Piper} بهترین کیفیت خروجی را ارائه کرد، هرچند
از دو مدل دیگر کندتر بود.

این نتایج نشان می‌دهند که برای یک چت‌بات صوتی محلی، بزرگ‌ترین
یا سریع‌ترین مدل الزاماً بهترین انتخاب نیست.

انتخاب مناسب باید بر اساس تعامل میان کیفیت، تأخیر، مصرف منابع
و عملکرد کل سامانه انجام شود.

در فصل بعد، نتیجه‌گیری کلی پروژه، محدودیت‌ها و مسیرهای توسعه
آینده ارائه خواهند شد.
